\documentclass{article}
\usepackage[T1,T2A]{fontenc}
\usepackage[utf8]{inputenc}
\usepackage[english,russian]{babel}
\usepackage{xcolor}
\usepackage{amsmath}
\usepackage{amsfonts,bm}
\usepackage{tikz}
\usetikzlibrary{patterns}
\usetikzlibrary{quotes,angles}
\usepackage{relsize}

\title{Решение Задачи №28}
\author{Серёжи Фуканчика\\34 Ю класс}
\date{Декабрь 2019}

\begin{document}

\maketitle

\section{Задача}
\begin{equation}
\sqrt{2}\sin^3{x}-\sin^2{x}+\sqrt{2}\sin{x}-1=0
\end{equation}

\section{Решение}
Попробую замену $t=\sin{x}$, получаю кубическое уравнение:
\begin{equation}
\sqrt{2}t^3-t^2+\sqrt{2}t-1=0
\end{equation}

Части $\sqrt{2}t^3-t^2$ и $\sqrt{2}t-1$ похожи друг на друга, так что попробую вынести $t^2$ за скобку:
\begin{equation}
t^2(\sqrt{2}t-1)+(\sqrt{2}t-1)=0
\end{equation}

Теперь можно вынести $\sqrt{2}t-1$ за скобку:
\begin{equation}
(\sqrt{2}t-1)(t^2+1)=0
\end{equation}

Это распадается на совокупность:
\begin{equation*}
\left[
\begin{array}{l}
\sqrt{2}t-1 = 0 \\
t^2+1 = 0
\end{array}
\right.
\end{equation*}

перенесу константы вправо:
\begin{equation*}
\left[
\begin{array}{l}
\sqrt{2}t = 1 \\
t^2 = -1
\end{array}
\right.
\end{equation*}

$t^2 = -1$ не имеет решения в школьной математике, так что его отбрасываю. Надо решить
$$
\sqrt{2}t = 1
$$
поделю обе стороны на $\sqrt{2}$:
$$
t = \frac{1}{\sqrt{2}}
$$
вернусь от замены к исходной переменной:
$$
\sin{x} = \frac{1}{\sqrt{2}}
$$
или в более знакомой форме:
$$
\sin{x} = \frac{\sqrt{2}}{2}
$$
А это табличное значение $\frac{\pi}{4}$ и $\frac{3\pi}{4}$.

\section{Ответ}
$x_1=\frac{\pi}{4}+2k\pi,k\in\mathbb{Z}$

$x_2=\frac{3\pi}{4}+2m\pi,m\in\mathbb{Z}$

\section{Проверка}

\end{document}
